\section*{Conclusion}

Holding operations and corporate assumptions fixed isolates the effect of gold-price model choice. Nested-candidate safeguards improve the current Bates and Hawkes fits. Historical option-surface OOS remains distinct and no original engine beats interpolated persistence. Numerical and terminal sensitivity prevent a corporate value ranking.

The copper extension admits three producing mines and adds a conditional five-year operating contribution. Historical snapshots support actual temporal testing, but the sparse rolling experiment does not beat IV persistence. Its positive fixed-origin Bates and Hawkes comparisons are exploratory. The combined output remains an attributable component-margin sensitivity, without a reconciled FCFF or equity bridge.

The persistent-cost assumption is prudential only relative to otherwise identical efficiency-improvement scenarios, not a lower bound. The contribution is an auditable interdisciplinary framework, not a certified value: commodity basis, probability measure, operating uncertainty, corporate accounting and finite-resource closure remain the priorities for economic validation.


The next required result is a validated physical commodity forecast for each admitted model, with an explicit link to option calibration and Barrick's realised prices. This is a prerequisite for the intended DCF, not an optional robustness check. Until price, operating and accounting gates are met, the current calculations remain conditional sensitivities.
